\documentclass[11pt]{article}
% jugeo-common.tex — shared preamble for all Judgment Geometry papers
% Include via: % jugeo-common.tex — shared preamble for all Judgment Geometry papers
% Include via: % jugeo-common.tex — shared preamble for all Judgment Geometry papers
% Include via: \input{jugeo-common}

% ── Packages ────────────────────────────────────────────────────────────
\usepackage{amsmath,amssymb,amsthm,mathtools}
\usepackage{tikz-cd}
\usepackage{listings}
\usepackage{xcolor}
\usepackage{xspace}
\usepackage{booktabs}
\usepackage{multirow}
\usepackage{enumitem}
\usepackage[expansion=false]{microtype}
\usepackage{hyperref}
\usepackage{cleveref}
% \usepackage{thmtools}  % disabled: not available in this TeX installation
\usepackage{float}
\usepackage{caption}
\usepackage{subcaption}
\usepackage{tabularx}
\usepackage{stmaryrd}       % \llbracket, \rrbracket
\usepackage{bbm}            % \mathbbm{1}

% ── Page geometry ───────────────────────────────────────────────────────
\usepackage[margin=1in]{geometry}

% ── Theorem environments ────────────────────────────────────────────────
\theoremstyle{plain}
\newtheorem{theorem}{Theorem}[section]
\newtheorem{lemma}[theorem]{Lemma}
\newtheorem{proposition}[theorem]{Proposition}
\newtheorem{corollary}[theorem]{Corollary}

\theoremstyle{definition}
\newtheorem{definition}[theorem]{Definition}
\newtheorem{example}[theorem]{Example}
\newtheorem{construction}[theorem]{Construction}

\theoremstyle{remark}
\newtheorem{remark}[theorem]{Remark}
\newtheorem{notation}[theorem]{Notation}

% ── Python listings style ──────────────────────────────────────────────
\definecolor{jugeoBlue}{HTML}{2563EB}
\definecolor{jugeoGreen}{HTML}{059669}
\definecolor{jugeoRed}{HTML}{DC2626}
\definecolor{jugeoPurple}{HTML}{7C3AED}
\definecolor{jugeoGray}{HTML}{6B7280}
\definecolor{jugeoBg}{HTML}{F8FAFC}

\lstdefinestyle{jugeo-python}{
  language=Python,
  basicstyle=\ttfamily\small,
  keywordstyle=\color{jugeoBlue}\bfseries,
  stringstyle=\color{jugeoGreen},
  commentstyle=\color{jugeoGray}\itshape,
  numberstyle=\tiny\color{jugeoGray},
  backgroundcolor=\color{jugeoBg},
  numbers=left,
  numbersep=6pt,
  frame=single,
  framerule=0.4pt,
  rulecolor=\color{jugeoGray!40},
  breaklines=true,
  breakatwhitespace=true,
  showstringspaces=false,
  tabsize=4,
  xleftmargin=1.5em,
  framexleftmargin=1.5em,
  morekeywords={dataclass,frozen,field,Enum,IntEnum,auto,Any,
                Mapping,Sequence,Optional,match,case},
  literate={→}{{$\to$}}1 {←}{{$\leftarrow$}}1
           {≤}{{$\leq$}}1 {≥}{{$\geq$}}1 {≠}{{$\neq$}}1
           {∈}{{$\in$}}1 {∀}{{$\forall$}}1 {∃}{{$\exists$}}1
           {⊕}{{$\oplus$}}1 {⊖}{{$\ominus$}}1,
  escapeinside={(*@}{@*)},
}
\lstset{style=jugeo-python}

\lstdefinestyle{jugeo-output}{
  basicstyle=\ttfamily\small,
  backgroundcolor=\color{jugeoBg},
  frame=single,
  framerule=0.4pt,
  rulecolor=\color{jugeoGray!40},
  breaklines=true,
  showstringspaces=false,
  xleftmargin=1.5em,
  framexleftmargin=0.5em,
  numbers=none,
}

% ── JuGeo-specific macros ──────────────────────────────────────────────

% Judgment 8-tuple
\newcommand{\J}{\mathcal{J}}
\newcommand{\Jud}[8]{(#1,\,#2,\,#3,\,#4,\,#5,\,#6,\,#7,\,#8)}
\newcommand{\JudShort}[1]{\J_{#1}}

% Components
\newcommand{\coord}{\mathsf{c}}            % coordinate
\newcommand{\prop}{\varphi}                 % proposition
\newcommand{\carrier}{\mathcal{A}}          % carrier type
\newcommand{\evid}{\mathcal{E}}             % evidence bundle
\newcommand{\oblig}{\mathcal{O}}            % obligations
\newcommand{\obstr}{\mathcal{B}}            % obstructions
\newcommand{\trust}{\mathcal{T}}            % trust annotation
\newcommand{\prov}{\Pi}                     % provenance

% Trust algebra
\newcommand{\Talg}{\mathfrak{T}}            % trust ordered algebra
\newcommand{\Eadm}{\mathcal{E}_{\mathrm{adm}}}
\newcommand{\tleq}{\preceq}                 % trust ordering
\newcommand{\tjoin}{\oplus}                 % conservative join
\newcommand{\tatten}{\ominus}               % attenuation
\newcommand{\tpromote}{\uparrow_\pi}        % promotion
\newcommand{\tdemote}{\downarrow_\chi}       % demotion

% Trust tiers
\newcommand{\Tcontra}{\ensuremath{\mathsf{CONTRADICTED}}}
\newcommand{\Tunver}{\ensuremath{\mathsf{UNVERIFIED}}}
\newcommand{\Tcopilot}{\ensuremath{\mathsf{COPILOT}}}
\newcommand{\Toracle}{\ensuremath{\mathsf{ORACLE}}}
\newcommand{\Truntime}{\ensuremath{\mathsf{RUNTIME}}}
\newcommand{\Tsolver}{\ensuremath{\mathsf{SOLVER}}}
\newcommand{\Tproof}{\ensuremath{\mathsf{PROOF}}}

% Site and geometry
\newcommand{\Site}{\mathcal{S}}             % semantic site
\newcommand{\Cov}{\mathrm{Cov}}            % covering family
\newcommand{\Ob}{\mathrm{Ob}}              % objects
\newcommand{\Mor}{\mathrm{Mor}}            % morphisms
\newcommand{\res}{\mathrm{res}}            % restriction
\newcommand{\Cech}{\check{C}}              % Čech complex
\newcommand{\Hcoh}[1]{H^{#1}}             % cohomology group

% Descent
\newcommand{\descent}{\mathrm{desc}}
\newcommand{\glue}{\mathrm{glue}}
\newcommand{\overlap}[2]{#1 \cap #2}

% Semantic moves
\newcommand{\VERIFY}{\textsc{Verify}}
\newcommand{\CONSTRUCT}{\textsc{Construct}}
\newcommand{\NORMALIZE}{\textsc{Normalize}}
\newcommand{\GLUE}{\textsc{Glue}}
\newcommand{\RESTRICT}{\textsc{Restrict}}
\newcommand{\TRANSPORT}{\textsc{Transport}}
\newcommand{\REPAIR}{\textsc{Repair}}
\newcommand{\PROMOTE}{\textsc{Promote}}
\newcommand{\CHALLENGE}{\textsc{Challenge}}
\newcommand{\ESCALATE}{\textsc{Escalate}}

% Fragments
\newcommand{\QFLIA}{\texttt{QF\_LIA}}
\newcommand{\QFLRA}{\texttt{QF\_LRA}}
\newcommand{\QFBV}{\texttt{QF\_BV}}
\newcommand{\QFUF}{\texttt{QF\_UF}}

% Misc
\newcommand{\Python}{\textsc{Python}}
\newcommand{\JuGeo}{\textsc{JuGeo}}
\newcommand{\LEAN}{\textsc{Lean}\,4}
\newcommand{\Fstar}{F$^{\star}$}
\newcommand{\Coq}{\textsc{Coq}}
\newcommand{\Dafny}{\textsc{Dafny}}

% Common shortcuts
\newcommand{\ie}{i.e.\@\xspace}
\newcommand{\eg}{e.g.\@\xspace}
\newcommand{\cf}{cf.\@\xspace}
\newcommand{\wrt}{w.r.t.\@\xspace}

% Evaluation shortcuts
\newcommand{\numcases}{300}
\newcommand{\numspec}{100}
\newcommand{\numeq}{100}
\newcommand{\numbug}{100}
\newcommand{\medtime}{6.5\,ms}
\newcommand{\totaltime}{1.949\,s}


% ── Packages ────────────────────────────────────────────────────────────
\usepackage{amsmath,amssymb,amsthm,mathtools}
\usepackage{tikz-cd}
\usepackage{listings}
\usepackage{xcolor}
\usepackage{xspace}
\usepackage{booktabs}
\usepackage{multirow}
\usepackage{enumitem}
\usepackage[expansion=false]{microtype}
\usepackage{hyperref}
\usepackage{cleveref}
% \usepackage{thmtools}  % disabled: not available in this TeX installation
\usepackage{float}
\usepackage{caption}
\usepackage{subcaption}
\usepackage{tabularx}
\usepackage{stmaryrd}       % \llbracket, \rrbracket
\usepackage{bbm}            % \mathbbm{1}

% ── Page geometry ───────────────────────────────────────────────────────
\usepackage[margin=1in]{geometry}

% ── Theorem environments ────────────────────────────────────────────────
\theoremstyle{plain}
\newtheorem{theorem}{Theorem}[section]
\newtheorem{lemma}[theorem]{Lemma}
\newtheorem{proposition}[theorem]{Proposition}
\newtheorem{corollary}[theorem]{Corollary}

\theoremstyle{definition}
\newtheorem{definition}[theorem]{Definition}
\newtheorem{example}[theorem]{Example}
\newtheorem{construction}[theorem]{Construction}

\theoremstyle{remark}
\newtheorem{remark}[theorem]{Remark}
\newtheorem{notation}[theorem]{Notation}

% ── Python listings style ──────────────────────────────────────────────
\definecolor{jugeoBlue}{HTML}{2563EB}
\definecolor{jugeoGreen}{HTML}{059669}
\definecolor{jugeoRed}{HTML}{DC2626}
\definecolor{jugeoPurple}{HTML}{7C3AED}
\definecolor{jugeoGray}{HTML}{6B7280}
\definecolor{jugeoBg}{HTML}{F8FAFC}

\lstdefinestyle{jugeo-python}{
  language=Python,
  basicstyle=\ttfamily\small,
  keywordstyle=\color{jugeoBlue}\bfseries,
  stringstyle=\color{jugeoGreen},
  commentstyle=\color{jugeoGray}\itshape,
  numberstyle=\tiny\color{jugeoGray},
  backgroundcolor=\color{jugeoBg},
  numbers=left,
  numbersep=6pt,
  frame=single,
  framerule=0.4pt,
  rulecolor=\color{jugeoGray!40},
  breaklines=true,
  breakatwhitespace=true,
  showstringspaces=false,
  tabsize=4,
  xleftmargin=1.5em,
  framexleftmargin=1.5em,
  morekeywords={dataclass,frozen,field,Enum,IntEnum,auto,Any,
                Mapping,Sequence,Optional,match,case},
  literate={→}{{$\to$}}1 {←}{{$\leftarrow$}}1
           {≤}{{$\leq$}}1 {≥}{{$\geq$}}1 {≠}{{$\neq$}}1
           {∈}{{$\in$}}1 {∀}{{$\forall$}}1 {∃}{{$\exists$}}1
           {⊕}{{$\oplus$}}1 {⊖}{{$\ominus$}}1,
  escapeinside={(*@}{@*)},
}
\lstset{style=jugeo-python}

\lstdefinestyle{jugeo-output}{
  basicstyle=\ttfamily\small,
  backgroundcolor=\color{jugeoBg},
  frame=single,
  framerule=0.4pt,
  rulecolor=\color{jugeoGray!40},
  breaklines=true,
  showstringspaces=false,
  xleftmargin=1.5em,
  framexleftmargin=0.5em,
  numbers=none,
}

% ── JuGeo-specific macros ──────────────────────────────────────────────

% Judgment 8-tuple
\newcommand{\J}{\mathcal{J}}
\newcommand{\Jud}[8]{(#1,\,#2,\,#3,\,#4,\,#5,\,#6,\,#7,\,#8)}
\newcommand{\JudShort}[1]{\J_{#1}}

% Components
\newcommand{\coord}{\mathsf{c}}            % coordinate
\newcommand{\prop}{\varphi}                 % proposition
\newcommand{\carrier}{\mathcal{A}}          % carrier type
\newcommand{\evid}{\mathcal{E}}             % evidence bundle
\newcommand{\oblig}{\mathcal{O}}            % obligations
\newcommand{\obstr}{\mathcal{B}}            % obstructions
\newcommand{\trust}{\mathcal{T}}            % trust annotation
\newcommand{\prov}{\Pi}                     % provenance

% Trust algebra
\newcommand{\Talg}{\mathfrak{T}}            % trust ordered algebra
\newcommand{\Eadm}{\mathcal{E}_{\mathrm{adm}}}
\newcommand{\tleq}{\preceq}                 % trust ordering
\newcommand{\tjoin}{\oplus}                 % conservative join
\newcommand{\tatten}{\ominus}               % attenuation
\newcommand{\tpromote}{\uparrow_\pi}        % promotion
\newcommand{\tdemote}{\downarrow_\chi}       % demotion

% Trust tiers
\newcommand{\Tcontra}{\ensuremath{\mathsf{CONTRADICTED}}}
\newcommand{\Tunver}{\ensuremath{\mathsf{UNVERIFIED}}}
\newcommand{\Tcopilot}{\ensuremath{\mathsf{COPILOT}}}
\newcommand{\Toracle}{\ensuremath{\mathsf{ORACLE}}}
\newcommand{\Truntime}{\ensuremath{\mathsf{RUNTIME}}}
\newcommand{\Tsolver}{\ensuremath{\mathsf{SOLVER}}}
\newcommand{\Tproof}{\ensuremath{\mathsf{PROOF}}}

% Site and geometry
\newcommand{\Site}{\mathcal{S}}             % semantic site
\newcommand{\Cov}{\mathrm{Cov}}            % covering family
\newcommand{\Ob}{\mathrm{Ob}}              % objects
\newcommand{\Mor}{\mathrm{Mor}}            % morphisms
\newcommand{\res}{\mathrm{res}}            % restriction
\newcommand{\Cech}{\check{C}}              % Čech complex
\newcommand{\Hcoh}[1]{H^{#1}}             % cohomology group

% Descent
\newcommand{\descent}{\mathrm{desc}}
\newcommand{\glue}{\mathrm{glue}}
\newcommand{\overlap}[2]{#1 \cap #2}

% Semantic moves
\newcommand{\VERIFY}{\textsc{Verify}}
\newcommand{\CONSTRUCT}{\textsc{Construct}}
\newcommand{\NORMALIZE}{\textsc{Normalize}}
\newcommand{\GLUE}{\textsc{Glue}}
\newcommand{\RESTRICT}{\textsc{Restrict}}
\newcommand{\TRANSPORT}{\textsc{Transport}}
\newcommand{\REPAIR}{\textsc{Repair}}
\newcommand{\PROMOTE}{\textsc{Promote}}
\newcommand{\CHALLENGE}{\textsc{Challenge}}
\newcommand{\ESCALATE}{\textsc{Escalate}}

% Fragments
\newcommand{\QFLIA}{\texttt{QF\_LIA}}
\newcommand{\QFLRA}{\texttt{QF\_LRA}}
\newcommand{\QFBV}{\texttt{QF\_BV}}
\newcommand{\QFUF}{\texttt{QF\_UF}}

% Misc
\newcommand{\Python}{\textsc{Python}}
\newcommand{\JuGeo}{\textsc{JuGeo}}
\newcommand{\LEAN}{\textsc{Lean}\,4}
\newcommand{\Fstar}{F$^{\star}$}
\newcommand{\Coq}{\textsc{Coq}}
\newcommand{\Dafny}{\textsc{Dafny}}

% Common shortcuts
\newcommand{\ie}{i.e.\@\xspace}
\newcommand{\eg}{e.g.\@\xspace}
\newcommand{\cf}{cf.\@\xspace}
\newcommand{\wrt}{w.r.t.\@\xspace}

% Evaluation shortcuts
\newcommand{\numcases}{300}
\newcommand{\numspec}{100}
\newcommand{\numeq}{100}
\newcommand{\numbug}{100}
\newcommand{\medtime}{6.5\,ms}
\newcommand{\totaltime}{1.949\,s}


% ── Packages ────────────────────────────────────────────────────────────
\usepackage{amsmath,amssymb,amsthm,mathtools}
\usepackage{tikz-cd}
\usepackage{listings}
\usepackage{xcolor}
\usepackage{xspace}
\usepackage{booktabs}
\usepackage{multirow}
\usepackage{enumitem}
\usepackage[expansion=false]{microtype}
\usepackage{hyperref}
\usepackage{cleveref}
% \usepackage{thmtools}  % disabled: not available in this TeX installation
\usepackage{float}
\usepackage{caption}
\usepackage{subcaption}
\usepackage{tabularx}
\usepackage{stmaryrd}       % \llbracket, \rrbracket
\usepackage{bbm}            % \mathbbm{1}

% ── Page geometry ───────────────────────────────────────────────────────
\usepackage[margin=1in]{geometry}

% ── Theorem environments ────────────────────────────────────────────────
\theoremstyle{plain}
\newtheorem{theorem}{Theorem}[section]
\newtheorem{lemma}[theorem]{Lemma}
\newtheorem{proposition}[theorem]{Proposition}
\newtheorem{corollary}[theorem]{Corollary}

\theoremstyle{definition}
\newtheorem{definition}[theorem]{Definition}
\newtheorem{example}[theorem]{Example}
\newtheorem{construction}[theorem]{Construction}

\theoremstyle{remark}
\newtheorem{remark}[theorem]{Remark}
\newtheorem{notation}[theorem]{Notation}

% ── Python listings style ──────────────────────────────────────────────
\definecolor{jugeoBlue}{HTML}{2563EB}
\definecolor{jugeoGreen}{HTML}{059669}
\definecolor{jugeoRed}{HTML}{DC2626}
\definecolor{jugeoPurple}{HTML}{7C3AED}
\definecolor{jugeoGray}{HTML}{6B7280}
\definecolor{jugeoBg}{HTML}{F8FAFC}

\lstdefinestyle{jugeo-python}{
  language=Python,
  basicstyle=\ttfamily\small,
  keywordstyle=\color{jugeoBlue}\bfseries,
  stringstyle=\color{jugeoGreen},
  commentstyle=\color{jugeoGray}\itshape,
  numberstyle=\tiny\color{jugeoGray},
  backgroundcolor=\color{jugeoBg},
  numbers=left,
  numbersep=6pt,
  frame=single,
  framerule=0.4pt,
  rulecolor=\color{jugeoGray!40},
  breaklines=true,
  breakatwhitespace=true,
  showstringspaces=false,
  tabsize=4,
  xleftmargin=1.5em,
  framexleftmargin=1.5em,
  morekeywords={dataclass,frozen,field,Enum,IntEnum,auto,Any,
                Mapping,Sequence,Optional,match,case},
  literate={→}{{$\to$}}1 {←}{{$\leftarrow$}}1
           {≤}{{$\leq$}}1 {≥}{{$\geq$}}1 {≠}{{$\neq$}}1
           {∈}{{$\in$}}1 {∀}{{$\forall$}}1 {∃}{{$\exists$}}1
           {⊕}{{$\oplus$}}1 {⊖}{{$\ominus$}}1,
  escapeinside={(*@}{@*)},
}
\lstset{style=jugeo-python}

\lstdefinestyle{jugeo-output}{
  basicstyle=\ttfamily\small,
  backgroundcolor=\color{jugeoBg},
  frame=single,
  framerule=0.4pt,
  rulecolor=\color{jugeoGray!40},
  breaklines=true,
  showstringspaces=false,
  xleftmargin=1.5em,
  framexleftmargin=0.5em,
  numbers=none,
}

% ── JuGeo-specific macros ──────────────────────────────────────────────

% Judgment 8-tuple
\newcommand{\J}{\mathcal{J}}
\newcommand{\Jud}[8]{(#1,\,#2,\,#3,\,#4,\,#5,\,#6,\,#7,\,#8)}
\newcommand{\JudShort}[1]{\J_{#1}}

% Components
\newcommand{\coord}{\mathsf{c}}            % coordinate
\newcommand{\prop}{\varphi}                 % proposition
\newcommand{\carrier}{\mathcal{A}}          % carrier type
\newcommand{\evid}{\mathcal{E}}             % evidence bundle
\newcommand{\oblig}{\mathcal{O}}            % obligations
\newcommand{\obstr}{\mathcal{B}}            % obstructions
\newcommand{\trust}{\mathcal{T}}            % trust annotation
\newcommand{\prov}{\Pi}                     % provenance

% Trust algebra
\newcommand{\Talg}{\mathfrak{T}}            % trust ordered algebra
\newcommand{\Eadm}{\mathcal{E}_{\mathrm{adm}}}
\newcommand{\tleq}{\preceq}                 % trust ordering
\newcommand{\tjoin}{\oplus}                 % conservative join
\newcommand{\tatten}{\ominus}               % attenuation
\newcommand{\tpromote}{\uparrow_\pi}        % promotion
\newcommand{\tdemote}{\downarrow_\chi}       % demotion

% Trust tiers
\newcommand{\Tcontra}{\ensuremath{\mathsf{CONTRADICTED}}}
\newcommand{\Tunver}{\ensuremath{\mathsf{UNVERIFIED}}}
\newcommand{\Tcopilot}{\ensuremath{\mathsf{COPILOT}}}
\newcommand{\Toracle}{\ensuremath{\mathsf{ORACLE}}}
\newcommand{\Truntime}{\ensuremath{\mathsf{RUNTIME}}}
\newcommand{\Tsolver}{\ensuremath{\mathsf{SOLVER}}}
\newcommand{\Tproof}{\ensuremath{\mathsf{PROOF}}}

% Site and geometry
\newcommand{\Site}{\mathcal{S}}             % semantic site
\newcommand{\Cov}{\mathrm{Cov}}            % covering family
\newcommand{\Ob}{\mathrm{Ob}}              % objects
\newcommand{\Mor}{\mathrm{Mor}}            % morphisms
\newcommand{\res}{\mathrm{res}}            % restriction
\newcommand{\Cech}{\check{C}}              % Čech complex
\newcommand{\Hcoh}[1]{H^{#1}}             % cohomology group

% Descent
\newcommand{\descent}{\mathrm{desc}}
\newcommand{\glue}{\mathrm{glue}}
\newcommand{\overlap}[2]{#1 \cap #2}

% Semantic moves
\newcommand{\VERIFY}{\textsc{Verify}}
\newcommand{\CONSTRUCT}{\textsc{Construct}}
\newcommand{\NORMALIZE}{\textsc{Normalize}}
\newcommand{\GLUE}{\textsc{Glue}}
\newcommand{\RESTRICT}{\textsc{Restrict}}
\newcommand{\TRANSPORT}{\textsc{Transport}}
\newcommand{\REPAIR}{\textsc{Repair}}
\newcommand{\PROMOTE}{\textsc{Promote}}
\newcommand{\CHALLENGE}{\textsc{Challenge}}
\newcommand{\ESCALATE}{\textsc{Escalate}}

% Fragments
\newcommand{\QFLIA}{\texttt{QF\_LIA}}
\newcommand{\QFLRA}{\texttt{QF\_LRA}}
\newcommand{\QFBV}{\texttt{QF\_BV}}
\newcommand{\QFUF}{\texttt{QF\_UF}}

% Misc
\newcommand{\Python}{\textsc{Python}}
\newcommand{\JuGeo}{\textsc{JuGeo}}
\newcommand{\LEAN}{\textsc{Lean}\,4}
\newcommand{\Fstar}{F$^{\star}$}
\newcommand{\Coq}{\textsc{Coq}}
\newcommand{\Dafny}{\textsc{Dafny}}

% Common shortcuts
\newcommand{\ie}{i.e.\@\xspace}
\newcommand{\eg}{e.g.\@\xspace}
\newcommand{\cf}{cf.\@\xspace}
\newcommand{\wrt}{w.r.t.\@\xspace}

% Evaluation shortcuts
\newcommand{\numcases}{300}
\newcommand{\numspec}{100}
\newcommand{\numeq}{100}
\newcommand{\numbug}{100}
\newcommand{\medtime}{6.5\,ms}
\newcommand{\totaltime}{1.949\,s}


\title{\textbf{Accounting for Trust When Machines Write Proofs}}

\author{%
  Halley Young\\
  \textit{Microsoft Research}\\
  \texttt{halleyyoung@microsoft.com}
}

\date{\today}

\begin{document}
\maketitle

\begin{abstract}
Modern verification pipelines combine evidence from radically
heterogeneous sources: SMT solvers return \texttt{unsat},
large-language models declare ``looks correct,'' runtime monitors
observe 95\% test-pass rates, and human reviewers attest
correctness by inspection.
No existing system tracks the \emph{provenance} of each piece
of evidence or prevents the silent upgrading of a weak claim
to a strong one.
We introduce the \emph{trust ordered algebra}
$\Talg = (\Eadm,\;\tleq,\;\tjoin,\;\tatten,\;\tpromote,\;\tdemote)$,
an eight-level bounded distributive lattice equipped with
six operations that govern how evidence composes, attenuates,
promotes (only with explicit justification), and demotes across
channels.
The algebra enforces three invariants:
\emph{no silent promotion},
\emph{conservative join}, and
\emph{challenge conservativity}.
We implement $\Talg$ in \Python{} as part of the \JuGeo{} system
and evaluate it on six aspects of the trust algebra---lattice
properties, promotion blocking, copilot self-promotion guards,
audit logging, trust profiles, and overhead---achieving
a silent-promotion block rate of~1.0.
A head-to-head comparison with \LEAN{}, \Fstar{}, \Dafny{},
and GPT-4o shows that \JuGeo{} is the only system that
tracks trust provenance, enforces channel ceilings,
and integrates AI-generated evidence without silent promotion.
\end{abstract}

\tableofcontents
\newpage

%% ═══════════════════════════════════════════════════════════════════════
\section{Introduction}\label{sec:intro}
%% ═══════════════════════════════════════════════════════════════════════

\subsection{The mixed-evidence problem}

Consider the following scenario, increasingly common in
practice.  A developer writes a merge-sort implementation
with three specifications---sortedness, permutation preservation,
and stability---and asks the tool-chain to verify all three.
The SMT solver Z3 discharges the sortedness lemma in 12\,ms
and returns \texttt{unsat}; a large-language model (LLM)
inspects the code and reports ``permutation preservation looks
correct''; a test suite of 10\,000 random inputs passes
at 95\%.  Should the developer trust the conjunction?

Every existing proof assistant answers this question in
a binary fashion.  In \LEAN{} and \Coq{}, a proposition
is either proved or it is not; there is no middle
ground~\cite{lean4,coq}.  \Fstar{} enriches the picture
with computation types and multi-monadic effects~\cite{fstar},
but still collapses all evidence into a binary verdict.
\Dafny{} delegates to Boogie and Z3~\cite{dafny}, and the
user sees only \texttt{Verified} or an error message.
In all cases, the system discards the provenance of each
piece of evidence: a mechanically checked Coq proof carries
the same \emph{type} as a solver-discharged SMT query, and
there is no record that the permutation claim rests on an
LLM's informal assertion rather than a formal proof.

This paper introduces a \emph{trust ordered algebra}
$\Talg = (\Eadm,\;\tleq,\;\tjoin,\;\tatten,\;\tpromote,\;\tdemote)$
that solves the mixed-evidence problem by making trust
a first-class algebraic object.  Instead of collapsing
evidence into ``proved'' or ``not proved,'' $\Talg$ assigns
every claim a position in an eight-level bounded distributive
lattice, tracks which \emph{channel} produced it, and enforces
three invariants that prevent the silent inflation of trust.

\subsection{Contributions}

\begin{enumerate}[leftmargin=*]
\item We define an eight-level trust lattice $(\Eadm, \tleq)$
      and prove it is a bounded distributive lattice
      (\Cref{sec:tiers}).
\item We equip the lattice with six operations forming
      the trust ordered algebra $\Talg$
      (\Cref{sec:algebra}) and state three invariants
      (\Cref{sec:invariants}).
\item We prove a soundness theorem: the trust level
      assigned to a judgment faithfully reflects the quality
      of its evidence (\Cref{sec:soundness}).
\item We formalize how AI-generated evidence enters the
      algebra and show it can never silently promote itself
      (\Cref{sec:ai}).
\item We evaluate $\Talg$ on six dimensions of trust-algebra
      behaviour---lattice properties, promotion blocking, copilot
      guards, audit logging, trust profiles, and a state-of-the-art
      comparison with four systems (\Cref{sec:eval}).
\item We develop trust profiles---immutable, composable
      summaries of clausewise trust---and show they form
      a semilattice (\Cref{sec:profiles}).
\end{enumerate}

\subsection{Notation}

We write $\Talg$ for the full algebra,
$\Eadm$ for its carrier set (admissible evidence configurations),
$\tleq$ for the partial order,
$\tjoin$ for the conservative join,
$\tatten$ for attenuation,
$\tpromote$ for (justified) promotion, and
$\tdemote$ for demotion.
Trust tiers are typeset in sans-serif:
$\Tcontra \tleq \Tunver \tleq \Tcopilot \tleq \Toracle
 \tleq \mathsf{HUMAN} \tleq \Truntime \tleq \Tsolver \tleq \Tproof$.
Individual judgments are denoted
$\J = (\coord, \prop, \carrier, \evid, \oblig, \obstr,
       \trust, \prov)$
following the 8-tuple of~\cite{paper02}.


%% ═══════════════════════════════════════════════════════════════════════
\section{Trust Tiers}\label{sec:tiers}
%% ═══════════════════════════════════════════════════════════════════════

\subsection{The eight-level lattice}

The carrier set $\Eadm$ consists of eight trust levels,
ordered from weakest to strongest:

\begin{definition}[Trust lattice]\label{def:trust-lattice}
Let $\Eadm = \{\Tcontra,\;\Tunver,\;\Tcopilot,\;\Toracle,
\;\mathsf{HUMAN},\;\Truntime,\;\Tsolver,\;\Tproof\}$ with the total order
\[
  \Tcontra \;\tleq\; \Tunver \;\tleq\; \Tcopilot
  \;\tleq\; \Toracle \;\tleq\; \mathsf{HUMAN}
  \;\tleq\; \Truntime
  \;\tleq\; \Tsolver \;\tleq\; \Tproof.
\]
Define $\mathsf{meet}(a,b) = \min(a,b)$ and
$\mathsf{join}(a,b) = \max(a,b)$ under this order.
Then $(\Eadm, \tleq, \mathsf{meet}, \mathsf{join},
       \Tcontra, \Tproof)$
is a bounded lattice with bottom $\bot = \Tcontra$ and
top $\top = \Tproof$.
\end{definition}

\noindent
The implementation instantiates this lattice as a \Python{} enum:

\begin{lstlisting}[caption={Trust levels (\texttt{src/jugeo/evidence/trust.py}).},
                   label=lst:trust-levels]
class TrustLevel(str, Enum):
    MECHANICALLY_VERIFIED = 'mechanically_verified'   # top
    SOLVER_DISCHARGED     = 'solver_discharged'
    RUNTIME_WITNESSED     = 'runtime_witnessed'
    HUMAN_ATTESTED        = 'human_attested'
    ORACLE_PROPOSED       = 'oracle_proposed'
    COPILOT_SUGGESTED     = 'copilot_suggested'
    UNVERIFIED            = 'unverified'
    CONTRADICTED          = 'contradicted'            # bottom
\end{lstlisting}

\noindent
\textbf{Remark.}
The formal definition and the implementation
(\Cref{lst:trust-levels}) both use eight levels; the
$\mathsf{HUMAN\_ATTESTED}$ tier sits between $\Toracle$
and $\Truntime$.  The algebra remains a bounded distributive
lattice over all eight levels.

\subsection{Hasse diagram}

\begin{figure}[t]
\centering
\begin{tikzcd}[row sep=1.4em, column sep=0em]
  & \Tproof \arrow[d, dash] & \\
  & \Tsolver \arrow[d, dash] & \\
  & \Truntime \arrow[d, dash] & \\
  & \mathsf{HUMAN} \arrow[d, dash] & \\
  & \Toracle \arrow[d, dash] & \\
  & \Tcopilot \arrow[d, dash] & \\
  & \Tunver \arrow[d, dash] & \\
  & \Tcontra &
\end{tikzcd}
\caption{Hasse diagram of the trust lattice $(\Eadm, \tleq)$.
  The lattice is totally ordered; every pair of elements is
  comparable.  Edges denote the covering relation.}
\label{fig:hasse-main}
\end{figure}

\subsection{The five-tier sub-lattice}

The hypercover treaty negotiation module uses a coarser,
five-level sub-lattice suitable for cross-module boundary
checking:

\begin{lstlisting}[caption={Five-tier lattice
  (\texttt{hypercover\_treaties/algorithms.py}).},
                   label=lst:five-tier]
class TrustTier(IntEnum):
    PROPOSAL         = 1
    REVIEWED         = 2
    VERIFIED         = 3
    RUNTIME_WITNESSED = 4
    PROOF_BACKED     = 5

    def join(self, other):
        return TrustTier(max(self.value, other.value))
    def meet(self, other):
        return TrustTier(min(self.value, other.value))
    def promote(self):
        return TrustTier(min(self.value + 1, 5))
    def demote(self):
        return TrustTier(max(self.value - 1, 1))
\end{lstlisting}

\begin{figure}[t]
\centering
\begin{tikzcd}[row sep=1.4em, column sep=0em]
  & \mathsf{PROOF\_BACKED} \arrow[d, dash] & \\
  & \mathsf{RUNTIME\_WITNESSED} \arrow[d, dash] & \\
  & \mathsf{VERIFIED} \arrow[d, dash] & \\
  & \mathsf{REVIEWED} \arrow[d, dash] & \\
  & \mathsf{PROPOSAL} &
\end{tikzcd}
\caption{Hasse diagram of the five-tier sub-lattice used in
  hypercover treaty negotiation.  The embedding
  $\iota \colon \mathsf{TrustTier} \hookrightarrow \Eadm$
  preserves meets and joins.}
\label{fig:hasse-five}
\end{figure}

\noindent
The two lattices are related by a lattice homomorphism
$\iota \colon \mathsf{TrustTier} \to \Eadm$ that maps
$\mathsf{PROPOSAL} \mapsto \Tunver$,
$\mathsf{REVIEWED} \mapsto \Toracle$,
$\mathsf{VERIFIED} \mapsto \Truntime$,
$\mathsf{RUNTIME\_WITNESSED} \mapsto \Tsolver$,
$\mathsf{PROOF\_BACKED} \mapsto \Tproof$.

\begin{proposition}[Embedding preserves lattice structure]
\label{prop:embedding}
The map $\iota$ is an injective lattice homomorphism:
for all $a, b \in \mathsf{TrustTier}$,
\[
  \iota(\mathsf{meet}(a,b)) = \mathsf{meet}(\iota(a), \iota(b))
  \qquad\text{and}\qquad
  \iota(\mathsf{join}(a,b)) = \mathsf{join}(\iota(a), \iota(b)).
\]
\end{proposition}

\begin{proof}
Both lattices are totally ordered and $\iota$ is
order-preserving (strictly monotone on the five-element chain).
On a total order, $\mathsf{meet} = \min$ and
$\mathsf{join} = \max$.
Since $\iota$ preserves the order,
$\iota(\min(a,b)) = \min(\iota(a), \iota(b))$ and
$\iota(\max(a,b)) = \max(\iota(a), \iota(b))$.
Injectivity is immediate from strict monotonicity.
\end{proof}

\subsection{Distributivity}

\begin{theorem}[Bounded distributive lattice]
\label{thm:distributive}
$(\Eadm, \tleq, \mathsf{meet}, \mathsf{join}, \Tcontra, \Tproof)$
is a bounded distributive lattice.
\end{theorem}

\begin{proof}
A totally ordered set (chain) is trivially a distributive lattice:
for any $a, b, c$ in a chain,
\begin{align*}
  a \wedge (b \vee c)
    &= \min(a, \max(b, c)) \\
    &= \max(\min(a, b), \min(a, c)) \\
    &= (a \wedge b) \vee (a \wedge c).
\end{align*}
The identity holds by case analysis on the three possible
orderings of $a$, $b$, $c$.
Boundedness holds by construction:
$\Tcontra \tleq x \tleq \Tproof$ for all $x \in \Eadm$.
\end{proof}

\begin{corollary}\label{cor:absorption}
The absorption laws hold: for all $a, b \in \Eadm$,
\[
  a \wedge (a \vee b) = a
  \qquad\text{and}\qquad
  a \vee (a \wedge b) = a.
\]
\end{corollary}


%% ═══════════════════════════════════════════════════════════════════════
\section{The Trust Algebra \texorpdfstring{$\Talg$}{T}}\label{sec:algebra}
%% ═══════════════════════════════════════════════════════════════════════

We now equip the bounded distributive lattice
$(\Eadm, \tleq)$ with six operations to form the full
trust ordered algebra.

\begin{definition}[Trust ordered algebra]\label{def:trust-algebra}
The \emph{trust ordered algebra} is the six-tuple
\[
  \Talg = (\Eadm,\;\tleq,\;\tjoin,\;\tatten,\;\tpromote,\;\tdemote)
\]
where:
\begin{enumerate}[label=(\roman*)]
\item $(\Eadm, \tleq)$ is the bounded distributive lattice
      of~\Cref{def:trust-lattice}.
\item $\tjoin \colon \Eadm \times \Eadm \to \Eadm$ is the
      \emph{conservative join} (defined below).
\item $\tatten \colon \Eadm \times \mathbb{N} \to \Eadm$ is
      \emph{attenuation}.
\item $\tpromote \colon \Eadm \times \mathsf{Just} \to \Eadm$
      is \emph{promotion} (requires a non-empty justification).
\item $\tdemote \colon \Eadm \times \Eadm \to \Eadm$ is
      \emph{demotion} (enforces a ceiling).
\end{enumerate}
We also define a derived \emph{challenge} operator
$\mathsf{ch} \colon \Eadm \to \Eadm$.
\end{definition}

\subsection{Conservative join (\texorpdfstring{$\tjoin$}{oplus})}

Unlike the lattice join $\vee = \max$, the \emph{conservative
join} $\tjoin$ is deliberately pessimistic: when combining
evidence from two channels, the result takes the \emph{weaker}
of the two.

\begin{definition}[Conservative join]\label{def:cons-join}
For $a, b \in \Eadm$, define
\[
  a \tjoin b \;=\; \min(a, b).
\]
\end{definition}

\noindent
The rationale is epistemic: when composing evidence from
heterogeneous sources, the combined claim is only as strong
as its weakest link.  If an SMT solver discharges a lemma
($\Tsolver$) but the permutation argument rests on an LLM
suggestion ($\Tcopilot$), the conjunction of the two claims
should carry trust $\Tcopilot$, not $\Tsolver$.

\begin{lstlisting}[caption={Conservative join in the trust algebra.},
                   label=lst:compose]
class TrustAlgebra:
    def compose(self, a: TrustLevel, b: TrustLevel) -> TrustLevel:
        """Conservative join: min under the trust ordering."""
        if a.rank_index() <= b.rank_index():
            return a
        return b

    def meet(self, a: TrustLevel, b: TrustLevel) -> TrustLevel:
        """Greatest lower bound (coincides with compose)."""
        return self.compose(a, b)

    def join(self, a: TrustLevel, b: TrustLevel) -> TrustLevel:
        """Least upper bound (max under ordering)."""
        if a.rank_index() >= b.rank_index():
            return a
        return b
\end{lstlisting}

\begin{proposition}[Associativity and commutativity of $\tjoin$]
\label{prop:join-assoc}
For all $a, b, c \in \Eadm$:
\begin{enumerate}[label=(\alph*)]
\item $a \tjoin b = b \tjoin a$ \quad (commutativity);
\item $(a \tjoin b) \tjoin c = a \tjoin (b \tjoin c)$
      \quad (associativity);
\item $a \tjoin a = a$ \quad (idempotence);
\item $\Tcontra \tjoin a = \Tcontra$ \quad (annihilation by $\bot$);
\item $\Tproof \tjoin a = a$ \quad (identity at $\top$).
\end{enumerate}
\end{proposition}

\begin{proof}
All properties follow from $\tjoin = \min$ on a total order.
(a)~$\min(a,b) = \min(b,a)$.
(b)~$\min(\min(a,b),c) = \min(a,b,c) = \min(a, \min(b,c))$.
(c)~$\min(a,a) = a$.
(d)~$\Tcontra$ is the bottom, so $\min(\Tcontra, a) = \Tcontra$.
(e)~$\Tproof$ is the top, so $\min(\Tproof, a) = a$.
\end{proof}

\subsection{Attenuation (\texorpdfstring{$\tatten$}{ominus})}

When evidence is \emph{transported} across module boundaries
or \emph{restricted} to a sub-specification, its trust
may be attenuated.  Attenuation decrements the trust level
by a non-negative integer number of steps.

\begin{definition}[Attenuation]\label{def:attenuation}
For $t \in \Eadm$ and $k \in \mathbb{N}$, define
\[
  t \;\tatten\; k \;=\; \max\bigl(\Tcontra,\;\mathsf{pred}^k(t)\bigr)
\]
where $\mathsf{pred}$ is the predecessor function on the chain
(with $\mathsf{pred}(\Tcontra) = \Tcontra$).
\end{definition}

\begin{lstlisting}[caption={Attenuation through transport hops.},
                   label=lst:attenuation]
class TrustAttenuation:
    def attenuate_through_transport(self, level, hops):
        """Decrement trust by one step per hop, floor at CONTRADICTED."""
        current = level
        for _ in range(hops):
            idx = _ORDERED_LEVELS.index(current)
            if idx == 0:
                break
            current = _ORDERED_LEVELS[idx - 1]
        return current

    def attenuation_distance(self, source, target):
        """Number of steps between source and target."""
        return abs(
            _ORDERED_LEVELS.index(source) -
            _ORDERED_LEVELS.index(target)
        )
\end{lstlisting}

\begin{proposition}[Attenuation is monotone and idempotent at zero]
\label{prop:atten-mono}
For all $t \in \Eadm$ and $k, k' \in \mathbb{N}$:
\begin{enumerate}[label=(\alph*)]
\item $t \tatten 0 = t$;
\item if $k \leq k'$, then $t \tatten k' \tleq t \tatten k$;
\item $t \tatten k \tleq t$;
\item $(t \tatten k) \tatten k' = t \tatten (k + k')$.
\end{enumerate}
\end{proposition}

\begin{proof}
(a)~$\mathsf{pred}^0(t) = t$.
(b)--(c)~each application of $\mathsf{pred}$ decreases or
preserves the rank; more applications yield a weakly lower result.
(d)~both sides apply $\mathsf{pred}$ a total of $k + k'$ times,
floored at $\Tcontra$.
\end{proof}

\subsection{Promotion (\texorpdfstring{$\tpromote$}{uparrow})}

Promotion is the only operation that can \emph{increase}
a trust level, and it requires an explicit, non-empty
justification.  This is the key mechanism for preventing
silent upgrades.

\begin{definition}[Promotion]\label{def:promotion}
Let $\mathsf{Just}$ be the set of non-empty justification
strings.  For $t \in \Eadm$ and $j \in \mathsf{Just}$, define
\[
  \tpromote(t, j) \;=\; \mathsf{succ}(t)
  \qquad\text{if } j \neq \varepsilon,
\]
where $\mathsf{succ}$ is the successor function
(with $\mathsf{succ}(\Tproof) = \Tproof$).
If $j = \varepsilon$, the operation is \emph{undefined}
(raises a promotion failure).
\end{definition}

\begin{lstlisting}[caption={Promotion requires justification.},
                   label=lst:promote]
class TrustPromotion:
    def apply_promotion(self, proposal):
        """Validate, apply, and record a promotion."""
        valid, reason = self.validate_promotion(proposal)
        if not valid:
            raise PromotionError(reason)
        self.record_promotion(proposal)
        return proposal.target

def promote(self, t, justification):
    if not justification or not justification.strip():
        raise _make_promotion_failure(t, 'empty justification')
    if t == TrustLevel.MECHANICALLY_VERIFIED:
        return t   # already at top
    idx = _ORDERED_LEVELS.index(t)
    return _ORDERED_LEVELS[idx + 1]
\end{lstlisting}

\noindent
Two additional policy constraints guard promotion:
\begin{itemize}
\item \textbf{Copilot ceiling.}
  A claim originating from a \texttt{copilot} channel cannot
  be promoted above $\Toracle$.
\item \textbf{Witness requirement.}
  Promotion above $\mathsf{HUMAN\_ATTESTED}$ requires an
  external witness (e.g., a solver certificate or a runtime
  trace).
\end{itemize}

\subsection{Demotion (\texorpdfstring{$\tdemote$}{downarrow})}

Demotion enforces a ceiling: it clamps a trust level down
to a specified maximum.

\begin{definition}[Demotion]\label{def:demotion}
For $t, c \in \Eadm$, define
\[
  \tdemote(t, c) \;=\; \min(t, c).
\]
\end{definition}

\noindent
Note that demotion and the conservative join coincide
algebraically ($\tdemote(t, c) = t \tjoin c$), but they
differ semantically: $\tjoin$ combines evidence from two
channels, while $\tdemote$ enforces an upper bound imposed
by a policy.

\subsection{Challenge}

\begin{definition}[Challenge]\label{def:challenge}
The \emph{challenge} operator decrements the trust level
by one step and narrows the scope of the associated profile:
\[
  \mathsf{ch}(t) \;=\; t \tatten 1.
\]
\end{definition}

\noindent
Challenges arise when contradictory evidence is discovered.
A failed test suite challenges a solver-discharged claim,
demoting it from $\Tsolver$ to $\Truntime$.

\subsection{Operator interaction diagram}

\begin{figure}[t]
\centering
\begin{tikzcd}[row sep=2.5em, column sep=3.5em]
  t_1 \arrow[r, "\tjoin"] \arrow[d, "\tatten\;k"']
    & t_1 \tjoin t_2 \arrow[d, "\tatten\;k"] \\
  t_1 \tatten k \arrow[r, "\tjoin"']
    & (t_1 \tjoin t_2) \tatten k
\end{tikzcd}
\caption{Interaction of conservative join and attenuation.
  The diagram commutes:
  $(t_1 \tjoin t_2) \tatten k = (t_1 \tatten k) \tjoin (t_2 \tatten k)$.}
\label{fig:interaction}
\end{figure}

\begin{proposition}[Attenuation distributes over $\tjoin$]
\label{prop:atten-distrib}
For all $t_1, t_2 \in \Eadm$ and $k \in \mathbb{N}$,
\[
  (t_1 \tjoin t_2) \tatten k
  \;=\;
  (t_1 \tatten k) \tjoin (t_2 \tatten k).
\]
\end{proposition}

\begin{proof}
Without loss of generality, assume $t_1 \tleq t_2$.
Then $t_1 \tjoin t_2 = t_1$, and the left-hand side is
$t_1 \tatten k$.
For the right-hand side,
$t_1 \tatten k \tleq t_2 \tatten k$
(since $\tatten k$ is monotone in $t$), so
$(t_1 \tatten k) \tjoin (t_2 \tatten k) = t_1 \tatten k$.
\end{proof}


%% ═══════════════════════════════════════════════════════════════════════
\section{Three Invariants}\label{sec:invariants}
%% ═══════════════════════════════════════════════════════════════════════

The trust algebra is governed by three invariants that
constrain how trust may evolve during verification.
These invariants are enforced both statically (by the
algebra's definitions) and dynamically (by the audit log).

\subsection{Invariant 1: No silent promotion}

\begin{definition}[No-silent-promotion invariant]
\label{def:nsp}
For every pair of consecutive trust states
$t_i, t_{i+1}$ in an execution trace, if
$t_i \tleq t_{i+1}$ with $t_i \neq t_{i+1}$
(strict increase), then the transition must be
a $\tpromote$ operation with a recorded, non-empty
justification~$j$.
\end{definition}

\begin{theorem}[No silent promotion]\label{thm:nsp}
Let $\sigma = t_0, t_1, \ldots, t_n$ be any execution
trace produced by the operations of $\Talg$.
Then for all $0 \leq i < n$:
\[
  t_i \;\tleq\; t_{i+1} \;\wedge\; t_i \neq t_{i+1}
  \quad\Longrightarrow\quad
  \exists\, j \in \mathsf{Just}.\;
  t_{i+1} = \tpromote(t_i, j).
\]
\end{theorem}

\begin{proof}
By exhaustive case analysis on the five operations of $\Talg$:
\begin{itemize}
\item $\tjoin = \min$: the result is $\tleq$ both inputs,
      hence $t_{i+1} \tleq t_i$.  A strict increase is
      impossible.
\item $\tatten$: $t \tatten k \tleq t$ for all $k \geq 0$
      (\Cref{prop:atten-mono}(c)).  Again, a strict increase
      is impossible.
\item $\tdemote(t, c) = \min(t, c) \tleq t$: cannot increase.
\item $\mathsf{ch}(t) = t \tatten 1 \tleq t$: cannot increase.
\item $\tpromote(t, j)$: the only operation that can yield
      $t_{i+1} \succ t_i$, and it requires $j \neq \varepsilon$
      by definition.
\end{itemize}
Hence every strict increase in the trace is a $\tpromote$
with a non-empty justification.
\end{proof}

\subsection{Invariant 2: Conservative join}

\begin{definition}[Conservative-join invariant]
\label{def:cj}
For all $a, b \in \Eadm$:
\[
  a \tjoin b \;\tleq\; a
  \qquad\text{and}\qquad
  a \tjoin b \;\tleq\; b.
\]
\end{definition}

\begin{proof}
Immediate from $\tjoin = \min$.
\end{proof}

\noindent
This invariant ensures that combining evidence from different
channels never inflates the trust level: the conjunction
of a solver proof and an LLM suggestion is trusted only
to the level of the LLM suggestion.

\subsection{Invariant 3: Challenge conservativity}

\begin{definition}[Challenge-conservativity invariant]
\label{def:cc}
For all $t \in \Eadm$:
\[
  \mathsf{ch}(t) \;\tleq\; t.
\]
\end{definition}

\begin{proof}
$\mathsf{ch}(t) = t \tatten 1 \tleq t$ by
\Cref{prop:atten-mono}(c).
\end{proof}

\noindent
Challenge conservativity guarantees that challenges only
weaken trust.  A challenged claim cannot emerge stronger
than before.

\begin{figure}[t]
\centering
\begin{tikzcd}[row sep=2em, column sep=4em]
  \Tsolver
    \arrow[r, "\tjoin\;\Tcopilot"]
    \arrow[d, "\mathsf{ch}"']
  & \Tcopilot
    \arrow[d, "\mathsf{ch}"] \\
  \Truntime
    \arrow[r, "\tjoin\;\Tcopilot"']
  & \Tcopilot
\end{tikzcd}
\caption{Invariant interaction: conservative join and
  challenge both commute and both monotonically decrease
  trust.  No path in the diagram increases trust.}
\label{fig:invariant-interaction}
\end{figure}


%% ═══════════════════════════════════════════════════════════════════════
\section{Soundness}\label{sec:soundness}
%% ═══════════════════════════════════════════════════════════════════════

We formalize a soundness theorem that connects the syntactic
trust assignment to the semantic quality of evidence.

\subsection{Evidence quality}

\begin{definition}[Evidence quality function]
\label{def:quality}
An \emph{evidence quality function}
$\mathsf{qual} \colon \mathsf{Evidence} \to \Eadm$
assigns to each piece of evidence a trust level according
to the following rules:
\begin{enumerate}[label=(\roman*)]
\item If $e$ is a machine-checked proof term accepted by
      a kernel (Coq, Lean, etc.), then
      $\mathsf{qual}(e) = \Tproof$.
\item If $e$ is an \texttt{unsat} certificate from an SMT
      solver, then $\mathsf{qual}(e) = \Tsolver$.
\item If $e$ is a runtime execution trace with passing
      assertions, then $\mathsf{qual}(e) = \Truntime$.
\item If $e$ is a human attestation, then
      $\mathsf{qual}(e) = \mathsf{HUMAN}$.
\item If $e$ is an oracle (external API) response, then
      $\mathsf{qual}(e) = \Toracle$.
\item If $e$ is an LLM-generated claim, then
      $\mathsf{qual}(e) = \Tcopilot$.
\item If $e$ is the absence of evidence, then
      $\mathsf{qual}(e) = \Tunver$.
\item If $e$ contradicts existing evidence, then
      $\mathsf{qual}(e) = \Tcontra$.
\end{enumerate}
\end{definition}

\subsection{The soundness theorem}

\begin{theorem}[Trust soundness]\label{thm:soundness}
Let $\J = (\coord, \prop, \carrier, \evid, \oblig, \obstr,
           \trust, \prov)$
be a judgment produced by the \JuGeo{} system.
Then:
\[
  \trust(\J) \;=\; \bigoplus_{e \in \evid(\J)} \mathsf{qual}(e)
\]
where $\bigoplus$ denotes the left-fold of $\tjoin$ over
the evidence bundle.
\end{theorem}

\begin{proof}[Proof sketch]
By induction on the derivation of $\J$.

\textbf{Base case.}
If $\J$ is a leaf judgment produced by a single evidence
channel $c$ with evidence $e$, then $\evid(\J) = \{e\}$
and $\trust(\J) = \mathsf{qual}(e)$ by construction
(each channel initializes trust at its quality level).

\textbf{Inductive case (composition).}
If $\J = \J_1 \tjoin \J_2$, then by the inductive hypothesis,
$\trust(\J_i) = \bigoplus_{e \in \evid(\J_i)} \mathsf{qual}(e)$
for $i = 1, 2$.
By the definition of judgment composition,
\begin{align*}
  \trust(\J)
    &= \trust(\J_1) \tjoin \trust(\J_2) \\
    &= \Bigl(\bigoplus_{e \in \evid(\J_1)} \mathsf{qual}(e)\Bigr)
       \tjoin
       \Bigl(\bigoplus_{e \in \evid(\J_2)} \mathsf{qual}(e)\Bigr) \\
    &= \bigoplus_{e \in \evid(\J_1) \cup \evid(\J_2)} \mathsf{qual}(e)
\end{align*}
where the last step uses associativity and commutativity
of $\tjoin$.

\textbf{Inductive case (attenuation).}
If $\J$ is obtained from $\J'$ by transport across $k$ hops,
then $\trust(\J) = \trust(\J') \tatten k$.
By~\Cref{prop:atten-distrib}, this equals
$\bigoplus_{e \in \evid(\J')} (\mathsf{qual}(e) \tatten k)$,
which is the composition of the attenuated qualities.

\textbf{Inductive case (promotion).}
If $\J$ is obtained from $\J'$ by promotion with
justification~$j$, then the audit log records the transition,
and the new trust level $\tpromote(\trust(\J'), j)$ is
consistent with the updated evidence bundle $\evid(\J') \cup \{j\}$.
\end{proof}

\subsection{Ceiling enforcement}

\begin{corollary}[Channel ceilings are sound]
\label{cor:ceilings}
If a judgment's evidence originates from channel $c$ with
ceiling $\gamma_c$, then $\trust(\J) \tleq \gamma_c$.
\end{corollary}

\begin{proof}
The \texttt{TrustCeiling} class enforces
$\tdemote(t, \gamma_c)$ on every trust assignment from
channel $c$.  Since $\tdemote(t, \gamma_c) = \min(t, \gamma_c)
\leq \gamma_c$, the bound holds.
\end{proof}

The default ceilings are:

\begin{center}
\begin{tabular}{ll}
\toprule
\textbf{Channel} & \textbf{Ceiling} \\
\midrule
\texttt{formal\_proof} & $\Tproof$ \\
\texttt{solver}        & $\Tsolver$ \\
\texttt{runtime}       & $\Truntime$ \\
\texttt{oracle}        & $\Toracle$ \\
\texttt{copilot}       & $\Tcopilot$ \\
\bottomrule
\end{tabular}
\end{center}


%% ═══════════════════════════════════════════════════════════════════════
\section{AI-Generated Evidence}\label{sec:ai}
%% ═══════════════════════════════════════════════════════════════════════

Large-language models (LLMs) are increasingly used to
generate code, specifications, and even informal ``proofs.''
The trust algebra provides a principled framework for
integrating such evidence without undermining the integrity
of the verification pipeline.

\subsection{Entry point}

All LLM-generated evidence enters the trust algebra at
level $\Tcopilot$.  This is the second-weakest level,
above only $\Tunver$ and $\Tcontra$.

\begin{definition}[AI evidence rule]\label{def:ai-entry}
For any evidence $e$ produced by an LLM:
\[
  \mathsf{qual}(e_{\mathsf{LLM}}) \;=\; \Tcopilot.
\]
\end{definition}

\subsection{Composition with stronger evidence}

When AI evidence is combined with evidence from a stronger
channel, the conservative join ensures the result is at
the AI level:

\begin{example}[Solver meets AI]\label{ex:solver-ai}
Let $e_1$ be a solver proof ($\Tsolver$) and $e_2$ be an
LLM claim ($\Tcopilot$).  Then:
\[
  \Tsolver \tjoin \Tcopilot
  \;=\; \min(\Tsolver, \Tcopilot)
  \;=\; \Tcopilot.
\]
The combined claim carries trust $\Tcopilot$.
The solver proof is not degraded---it retains its own trust
level---but the \emph{conjunction} is trusted only to the
level of its weakest component.
\end{example}

\noindent
This is exactly the behavior required for the merge-sort
scenario of~\Cref{sec:intro}: the sortedness lemma
(solver-discharged) and the permutation claim (LLM-suggested)
together carry trust $\Tcopilot$.

\subsection{Promotion path for AI evidence}

An LLM claim can be promoted, but only through explicit,
justified steps:

\begin{enumerate}[leftmargin=*]
\item $\Tcopilot \xrightarrow{\tpromote(\text{``human reviewed''})}
      \Toracle$: a human expert reviews the LLM output and
      attests its plausibility.
\item $\Toracle \xrightarrow{\tpromote(\text{``test suite passes''})}
      \Truntime$: a comprehensive test suite validates the claim.
\item $\Truntime \xrightarrow{\tpromote(\text{``SMT discharged''})}
      \Tsolver$: an SMT solver formally discharges the property.
\item $\Tsolver \xrightarrow{\tpromote(\text{``Lean proof''})}
      \Tproof$: a machine-checked proof is constructed.
\end{enumerate}

\noindent
Each step requires a non-empty justification string and is
recorded in the audit log.  The copilot ceiling (\Cref{cor:ceilings})
prevents an LLM from self-promoting above $\Toracle$;
further promotion requires evidence from a different channel.

\begin{figure}[t]
\centering
\begin{tikzcd}[row sep=2em, column sep=2.8em]
  \Tcopilot
    \arrow[r, "\tpromote"]
    \arrow[dr, dashed, "\text{blocked}"']
  & \Toracle
    \arrow[r, "\tpromote"]
  & \Truntime
    \arrow[r, "\tpromote"]
  & \Tsolver
    \arrow[r, "\tpromote"]
  & \Tproof \\
  & \Tsolver
    \arrow[u, crossing over, leftarrow, "\tjoin"']
  & & &
\end{tikzcd}
\caption{Promotion path for AI-generated evidence.
  The dashed arrow shows the blocked self-promotion:
  the copilot ceiling prevents direct escalation to
  $\Tsolver$.  The cross-over arrow shows how a solver
  proof, when joined with a copilot claim, yields
  $\Tcopilot$ (conservative join).}
\label{fig:ai-promotion}
\end{figure}

\subsection{Self-promotion guard}

\begin{proposition}[AI cannot self-promote]
\label{prop:no-self-promote}
Let $e$ be evidence produced by an LLM channel.
Then there is no sequence of operations within the
LLM channel alone that raises $\mathsf{qual}(e)$
above $\Toracle$.
\end{proposition}

\begin{proof}
The copilot ceiling $\gamma_{\mathsf{copilot}} = \Tcopilot$
ensures $\tdemote(t, \gamma_{\mathsf{copilot}}) \tleq \Tcopilot$
for any $t$.  Promotion within the copilot channel is
policy-blocked: the \texttt{copilot\_cannot\_self\_promote}
check rejects any $\tpromote$ originating from the copilot
channel when the target exceeds $\Toracle$.
By the no-silent-promotion invariant (\Cref{thm:nsp}),
no other operation can increase trust.
\end{proof}


%% ═══════════════════════════════════════════════════════════════════════
\section{Evaluation}\label{sec:eval}
%% ═══════════════════════════════════════════════════════════════════════

\subsection{Experimental setup}

We evaluate six aspects of the trust algebra empirically:
lattice properties (\S\ref{sec:eval-lattice}),
promotion blocking (\S\ref{sec:eval-promo}),
copilot self-promotion guards (\S\ref{sec:eval-copilot}),
audit log analysis (\S\ref{sec:eval-audit}),
trust profiles (\S\ref{sec:eval-profiles}),
and a state-of-the-art comparison (\S\ref{sec:eval-sota}).
All experiments are executed in the \JuGeo{} \Python{} implementation
and the numerical results are drawn from the companion data file
\texttt{results\_paper04.json}.

%%--------------------------------------------------------------------
\subsection{Lattice properties}\label{sec:eval-lattice}

\begin{table}[t]
\centering
\caption{Lattice join (least upper bound) and meet (greatest lower bound) samples.
All 64 composition pairs were tested; 8 representative samples are shown.}
\label{tab:lattice-join-meet}
\small
\begin{tabular}{@{}llclll@{}}
\toprule
$a$ & $b$ & & $a \vee b$ (join) & & $a \wedge b$ (meet) \\
\midrule
$\Tcontra$ & $\Tunver$ && $\Tunver$ && $\Tcontra$ \\
$\Tcontra$ & $\Toracle$ && $\Toracle$ && $\Tcontra$ \\
$\Tcontra$ & $\Truntime$ && $\Truntime$ && $\Tcontra$ \\
$\Tcontra$ & $\Tproof$ && $\Tproof$ && $\Tcontra$ \\
$\Tcopilot$ & $\Tunver$ && $\Tcopilot$ && $\Tunver$ \\
$\Tcopilot$ & $\Toracle$ && $\Toracle$ && $\Tcopilot$ \\
$\Tcopilot$ & $\Truntime$ && $\Truntime$ && $\Tcopilot$ \\
$\Tcopilot$ & $\Tproof$ && $\Tproof$ && $\Tcopilot$ \\
\bottomrule
\end{tabular}
\end{table}

\noindent
All 64 pairwise join and meet operations were verified to be
consistent with the total order on~$\Eadm$.
The promote chain traverses all 8 levels in 7 steps:
\[
  \Tcontra \to \Tunver \to \Tcopilot \to \Toracle
  \to \mathsf{HUMAN} \to \Truntime \to \Tsolver \to \Tproof.
\]

\begin{table}[t]
\centering
\caption{Demotion with ceiling $\Truntime$: $\tdemote(t, \Truntime) = \min(t, \Truntime)$.}
\label{tab:demote-ceiling}
\begin{tabular}{@{}ll@{}}
\toprule
\textbf{Input level} $t$ & \textbf{Result} $\tdemote(t, \Truntime)$ \\
\midrule
$\Tcontra$ & $\Tcontra$ \\
$\Tunver$ & $\Tunver$ \\
$\Tcopilot$ & $\Tcopilot$ \\
$\Toracle$ & $\Toracle$ \\
$\mathsf{HUMAN}$ & $\mathsf{HUMAN}$ \\
$\Truntime$ & $\Truntime$ \\
$\Tsolver$ & $\Truntime$ \\
$\Tproof$ & $\Truntime$ \\
\bottomrule
\end{tabular}
\end{table}

\noindent
\Cref{tab:demote-ceiling} confirms that demotion with
ceiling $\Truntime$ caps every level above $\Truntime$
to $\Truntime$ while leaving lower levels unchanged.

%%--------------------------------------------------------------------
\subsection{Promotion blocking}\label{sec:eval-promo}

\begin{table}[t]
\centering
\caption{Promotion blocking: justified promotions vs.\ empty-justification attempts.}
\label{tab:promotion-blocking}
\begin{tabular}{@{}lccc@{}}
\toprule
\textbf{Scenario} & \textbf{Accepted} & \textbf{Blocked} & \textbf{Block rate} \\
\midrule
Justified promotion (7 steps) & 7 & 0 & 0\% \\
Empty justification (7 steps) & 0 & 7 & 100\% \\
\midrule
\textbf{Silent promotion block rate} & \multicolumn{3}{c}{\textbf{1.0}} \\
\bottomrule
\end{tabular}
\end{table}

\noindent
All 7 justified promotions along the full chain from
$\Tcontra$ to $\Tproof$ succeed.
All 7 empty-justification attempts are blocked, yielding
a silent promotion block rate of~1.0 (\Cref{tab:promotion-blocking}).

%%--------------------------------------------------------------------
\subsection{Copilot self-promotion guard}\label{sec:eval-copilot}

\begin{table}[t]
\centering
\caption{Copilot self-promotion guard: attempts from the $\Tcopilot$ channel to reach each target level.}
\label{tab:copilot-guard}
\begin{tabular}{@{}lcc@{}}
\toprule
\textbf{Target level} & \textbf{Blocked?} & \textbf{Reason} \\
\midrule
$\Tcontra$ & No & Demotion, not promotion \\
$\Tunver$ & No & Demotion, not promotion \\
$\Tcopilot$ & No & Identity, not promotion \\
$\Toracle$ & No & Within copilot ceiling \\
$\mathsf{HUMAN}$ & \textbf{Yes} & Above copilot ceiling \\
$\Truntime$ & \textbf{Yes} & Above copilot ceiling \\
$\Tsolver$ & \textbf{Yes} & Above copilot ceiling \\
$\Tproof$ & \textbf{Yes} & Above copilot ceiling \\
\bottomrule
\end{tabular}
\end{table}

\noindent
Of 8 attempted self-promotions from the copilot channel,
4 are blocked (targets above $\Toracle$) and 4 are allowed
(targets at or below the copilot ceiling).
The GPT-4o scenario---where an LLM assigned at $\Tcopilot$
attempts to reach $\Tsolver$---is blocked with reason:
``copilot channel cannot promote above ORACLE\_PROPOSED.''

%%--------------------------------------------------------------------
\subsection{Audit log}\label{sec:eval-audit}

\begin{table}[t]
\centering
\caption{Audit log analysis on a module with injected silent promotions.}
\label{tab:audit}
\begin{tabular}{@{}lr@{}}
\toprule
\textbf{Metric} & \textbf{Value} \\
\midrule
Total audit entries & 9 \\
Silent promotions detected & 6 \\
Anomalies flagged & 1 \\
Silent promotion coordinates & \texttt{module.risky.fn\{0..4\}}, \texttt{module.suspicious.escalation} \\
\bottomrule
\end{tabular}
\end{table}

\noindent
The audit log captures 9 entries, identifies 6 silent promotions
at specific semantic coordinates, and flags 1 anomaly
(\Cref{tab:audit}).
All injected silent promotions are detected (100\% recall).

%%--------------------------------------------------------------------
\subsection{Trust profiles}\label{sec:eval-profiles}

\begin{table}[t]
\centering
\caption{Trust profiles and their pairwise conservative joins.}
\label{tab:profiles}
\begin{tabular}{@{}llll@{}}
\toprule
\textbf{Entity} & \textbf{Tier} & \textbf{Scope} & \textbf{Joined tier} \\
\midrule
engine\_0 & PROPOSAL & arithmetic & --- \\
engine\_1 & REVIEWED & heap, identity & --- \\
engine\_2 & VERIFIED & semantic & --- \\
\midrule
engine\_0 $\tjoin$ engine\_1 & --- & --- & PROPOSAL \\
engine\_0 $\tjoin$ engine\_2 & --- & --- & PROPOSAL \\
engine\_1 $\tjoin$ engine\_2 & --- & --- & REVIEWED \\
\midrule
\textbf{All three joined} & --- & --- & \textbf{PROPOSAL} \\
\bottomrule
\end{tabular}
\end{table}

\noindent
Profile composition confirms that the conservative join takes
the minimum tier (\Cref{tab:profiles}).
Joining all three profiles yields PROPOSAL---the weakest tier
among them.
The demote example shows VERIFIED can be demoted to PROPOSAL
(2~steps).

%%--------------------------------------------------------------------
\subsection{State-of-the-art comparison}\label{sec:eval-sota}

\begin{table}[t]
\centering
\caption{State-of-the-art comparison of trust-management capabilities.}
\label{tab:sota}
\small
\begin{tabular}{@{}lccccc@{}}
\toprule
\textbf{Feature}
  & \textbf{\LEAN{}}
  & \textbf{\Fstar{}}
  & \textbf{\Dafny{}}
  & \textbf{GPT-4o}
  & \textbf{\JuGeo{}} \\
\midrule
Trust levels
  & 1 & 2 & 1 & 0 & \textbf{8} \\
Trust lattice
  & --- & --- & --- & --- & \checkmark \\
Audit log
  & --- & --- & --- & --- & \checkmark \\
Copilot ceiling
  & N/A & N/A & N/A & N/A & $\Tcopilot$ \\
Silent promotion possible
  & Yes & Yes & Yes & Yes & \textbf{No} \\
\bottomrule
\end{tabular}
\end{table}

\noindent
\LEAN{} operates at a single trust level---a kernel-checked
proof and a \texttt{sorry} are structurally identical
post-acceptance~\cite{lean4}.
\Fstar{} distinguishes computational purity via effect
annotations (Tot, Dv, ST) but provides no trust ordering
or provenance tracking~\cite{fstar}.
\Dafny{} delegates all verification conditions to Boogie/Z3
with no trust metadata~\cite{dafny}.
GPT-4o generates claims with no formal guarantee whatsoever;
\JuGeo{} assigns such evidence to $\Tcopilot$ with no
silent promotion possible.
\JuGeo{} is the only system providing all five capabilities:
an 8-level lattice, algebraic composition, an append-only
audit log, copilot channel ceilings, and the
no-silent-promotion invariant.

%%--------------------------------------------------------------------
\subsection{Overhead analysis}

The trust algebra adds negligible overhead.
All experiments complete in a total elapsed time of 2.5\,ms.
Trust-level computation (comparison, composition, ceiling
enforcement, and audit logging) is dominated by the overhead
of the \Python{} runtime rather than the algebraic operations
themselves.


%% ═══════════════════════════════════════════════════════════════════════
\section{Trust Profiles}\label{sec:profiles}
%% ═══════════════════════════════════════════════════════════════════════

Individual trust levels track the provenance of single
evidence claims.  In practice, a verification result
spans multiple specifications and multiple coordinates.
\emph{Trust profiles} aggregate this information into
an immutable, composable summary.

\subsection{Definition}

\begin{definition}[Trust profile]\label{def:profile}
A \emph{trust profile} is a tuple
$P = (\mathsf{id},\;\tau,\;\mathsf{scope},\;\mathsf{reasons})$
where:
\begin{itemize}
\item $\mathsf{id}$ is the entity identifier;
\item $\tau \in \Eadm$ is the trust tier;
\item $\mathsf{scope} \subseteq \mathsf{Labels}$ is the set
      of specification labels covered;
\item $\mathsf{reasons}$ is an ordered tuple of justification
      strings.
\end{itemize}
\end{definition}

\begin{lstlisting}[caption={Trust profile implementation.},
                   label=lst:profile]
@dataclass(frozen=True)
class TrustProfile:
    entity_id: str
    tier: TrustTier
    support_scope: tuple[str, ...]
    reasons: tuple[str, ...]

    def join(self, other):
        """Conservative join: weaker tier, intersected scope."""
        new_tier = TrustTier(min(self.tier.value,
                                 other.tier.value))
        common = tuple(s for s in self.support_scope
                       if s in other.support_scope)
        scope = common if common else tuple(sorted(set(
            self.support_scope + other.support_scope)))
        return TrustProfile(
            entity_id=self.entity_id,
            tier=new_tier,
            support_scope=scope,
            reasons=self.reasons + other.reasons,
        )
\end{lstlisting}

\subsection{Profile composition}

The join of two profiles uses the conservative join on
tiers and intersects the support scopes (falling back
to union when the intersection is empty).

\begin{proposition}[Profile join is a semilattice operation]
\label{prop:profile-semilattice}
The set of trust profiles over a fixed entity, equipped
with the join of \Cref{lst:profile}, forms a
join-semilattice.
\end{proposition}

\begin{proof}
We verify the semilattice axioms.

\textbf{Commutativity.}
$P_1.\mathsf{join}(P_2) = P_2.\mathsf{join}(P_1)$:
the tier is $\min(\tau_1, \tau_2) = \min(\tau_2, \tau_1)$,
the scope intersection is symmetric, and the reasons tuple
is concatenated (so the order may differ, but the set of
reasons is the same; we consider profiles equivalent up to
reason reordering).

\textbf{Associativity.}
$(P_1.\mathsf{join}(P_2)).\mathsf{join}(P_3)
= P_1.\mathsf{join}(P_2.\mathsf{join}(P_3))$:
the tier is $\min(\tau_1, \tau_2, \tau_3)$ in both cases;
the scope is the intersection of all three (or the fallback
union in both cases, since the fallback triggers identically).

\textbf{Idempotence.}
$P.\mathsf{join}(P) = P$: $\min(\tau, \tau) = \tau$,
the scope intersection with itself is itself, and the
reasons are duplicated (but semantically equivalent).
\end{proof}

\subsection{Propagation across modules}

When a trust profile crosses a module boundary during
hypercover treaty negotiation, it is attenuated:

\begin{lstlisting}[caption={Profile propagation with attenuation.},
                   label=lst:propagation]
def propagate_profile(profile, boundary_hops):
    """Attenuate trust when crossing module boundaries."""
    new_tier = profile.tier
    for _ in range(boundary_hops):
        new_tier = new_tier.demote()
    return TrustProfile(
        entity_id=profile.entity_id,
        tier=new_tier,
        support_scope=profile.support_scope,
        reasons=profile.reasons + (
            f'attenuated-{boundary_hops}-hops',),
    )
\end{lstlisting}

\begin{proposition}[Propagation preserves conservativity]
\label{prop:propagation}
For any profile $P$ and hop count $k \geq 0$:
\[
  \mathsf{propagate}(P, k).\tau \;\tleq\; P.\tau.
\]
\end{proposition}

\begin{proof}
Each call to $\mathsf{demote}$ decreases the tier
by one step (floored at $\mathsf{PROPOSAL}$).
After $k$ demotions, the resulting tier is
$\tleq$ the original.
\end{proof}


%% ═══════════════════════════════════════════════════════════════════════
\section{Related Work and Conclusion}\label{sec:related}
%% ═══════════════════════════════════════════════════════════════════════

\subsection{Related work}

\paragraph{Trusted computing bases.}
The concept of a trusted computing base (TCB)~\cite{sel4,certikos}
identifies the minimal set of components that must be
correct for the system's security guarantees to hold.
Our trust algebra refines this idea: instead of a binary
``in the TCB / not in the TCB'' distinction, the eight-level
lattice tracks the degree of trust with fine granularity.
The no-silent-promotion invariant ensures that the TCB
boundary is not silently crossed.

\paragraph{Confidence logics.}
Probabilistic and possibilistic logics assign numerical
confidence scores to propositions~\cite{halpern03}.
Shafer's theory of evidence~\cite{shafer76} provides
belief functions on a lattice of hypotheses.
These approaches differ from ours in two respects:
(i)~they use continuous scores rather than a discrete
lattice, making auditing more difficult;
(ii)~they lack the operational structure (promotion,
demotion, channel ceilings) that our algebra provides.

\paragraph{Provenance semirings.}
Green, Karvounarakis, and Tannen~\cite{green07} introduced
provenance semirings for database query evaluation, tracking
\emph{how} a tuple was derived.  Buneman et~al.~\cite{buneman01}
developed the why/where/how provenance taxonomy.
Our trust algebra can be seen as a provenance semiring
specialized to verification evidence, with the conservative
join playing the role of the semiring multiplication.

\paragraph{Proof assistants.}
\LEAN{}~\cite{lean4}, \Coq{}~\cite{coq}, \Fstar{}~\cite{fstar},
and \Dafny{}~\cite{dafny} are the closest comparators.
As detailed in~\Cref{tab:sota}, all operate with at most
two implicit trust levels and none tracks evidence provenance
or enforces channel ceilings.
Isabelle/HOL~\cite{isabelle} uses a small kernel and
untrusted tactics, providing an implicit two-level trust
distinction, but this is not exposed to the user.

\paragraph{Refinement types and liquid types.}
Liquid Haskell~\cite{liquidhaskell} uses SMT-backed refinement
types to verify Haskell programs.  Like \Dafny{}, it delegates
to Z3 and presents a binary verdict.
The trust algebra could be layered on top of such systems
to track the provenance of each discharged refinement.

\paragraph{Separation logic.}
Reynolds's separation logic~\cite{reynolds02} and its
modern descendants (Iris~\cite{iris}) provide powerful
reasoning principles for shared mutable state.
These logics are orthogonal to our contribution: they
determine \emph{what} can be proved, while the trust
algebra tracks \emph{how much} each proof should be
trusted.

\subsection{An adversarial model for trust escalation}

We formalize the trust algebra's security properties using
a game-theoretic model.  Consider an adversary $\mathcal{A}$
(modeling a compromised LLM or malicious oracle) that seeks to
escalate a judgment's trust to $\Tsolver$ or $\Tproof$ without
providing genuine evidence at those levels.

\begin{definition}[Trust escalation game]\label{def:trust-game}
The \emph{trust escalation game}
$\mathcal{G} = (\mathcal{A}, \mathcal{V}, \Talg, k)$ is a
two-player game between an adversary $\mathcal{A}$ and a
verifier $\mathcal{V}$ over $k$ rounds:
\begin{enumerate}[nosep,label=(\roman*)]
  \item $\mathcal{A}$ starts with a judgment at trust level
        $\Tcopilot$ and may perform: compositions ($\tjoin$),
        restrictions, transports, and promotion requests.
  \item $\mathcal{V}$ enforces the trust algebra: each promotion
        request is checked against the policy route; compositions
        use conservative join.
  \item $\mathcal{A}$ wins if after $k$ rounds the judgment reaches
        $\Tsolver$ or above.
\end{enumerate}
\end{definition}

\begin{theorem}[Adversarial trust bound]\label{thm:adversarial}
No adversary $\mathcal{A}$ can win the trust escalation game
$\mathcal{G}$, regardless of $k$, provided the verifier $\mathcal{V}$
faithfully implements $\Talg$.
\end{theorem}

\begin{proof}
We prove by induction on $k$ that the maximum achievable trust
level is $\Toracle$ (oracle-proposed).

\emph{Base case} ($k = 0$): the judgment is at $\Tcopilot$.

\emph{Inductive step}: assume after $k$ rounds the maximum
achievable level is $\ell \tleq \Toracle$.  In round $k+1$,
$\mathcal{A}$ may:
\begin{itemize}[nosep]
  \item \emph{Compose}: $\ell \tjoin \ell' = \min(\ell, \ell')
        \tleq \ell$ (conservative join cannot increase trust).
  \item \emph{Restrict/transport}: by trust monotonicity
        (Theorem~\ref{thm:trust-monotone}), $\ell' \tleq \ell$.
  \item \emph{Request promotion}: the copilot channel ceiling
        (Theorem~\ref{thm:copilot-ceiling}) blocks promotion above
        $\Toracle$.  To reach $\Tsolver$, $\mathcal{A}$ needs
        evidence from a solver channel, which it does not possess
        (by assumption: $\mathcal{A}$ controls only copilot/oracle
        channels).
\end{itemize}
In all cases, $\ell' \tleq \Toracle < \Tsolver$, so $\mathcal{A}$
cannot reach $\Tsolver$.
\end{proof}

\begin{theorem}[Impossibility of silent promotion in flat systems]
\label{thm:flat-impossible}
Any verification system with a flat trust model (at most 2 trust
levels) cannot distinguish between the following two scenarios:
\begin{enumerate}[nosep,label=(\alph*)]
  \item A judgment fully proved by a trusted SMT solver.
  \item A judgment proposed by an LLM with no independent verification.
\end{enumerate}
In contrast, $\Talg$ assigns scenario (a) to $\Tsolver$ and
scenario (b) to $\Tcopilot$, with composition $\Tsolver \tjoin
\Tcopilot = \Tcopilot$ ensuring the LLM evidence cannot ``launder''
to solver-level trust.
\end{theorem}

\begin{proof}
In a flat system with trust levels $\{0, 1\}$, both scenarios
result in trust $1$ (``verified'').  There is no intermediate level
to which LLM evidence can be assigned while still being marked as
evidence (assigning $0$ would discard it entirely).  The composition
$1 \oplus 1 = 1$ cannot decrease, so any evidence at level $1$
permanently contaminates the trust pool.  In $\Talg$, the 8-level
lattice provides the necessary room: $\Tcopilot$ is strictly below
$\Tsolver$, and conservative join ($\min$) propagates the lower level.
\end{proof}

\subsection{Conclusion}

We have introduced the trust ordered algebra
$\Talg = (\Eadm, \tleq, \tjoin, \tatten, \tpromote, \tdemote)$,
an eight-level bounded distributive lattice with six operations
that governs how evidence composes, attenuates, and promotes
in a mixed-evidence verification pipeline.
The algebra enforces three invariants---no silent promotion,
conservative join, and challenge conservativity---and we
have proved soundness: the trust level of a judgment
faithfully reflects the quality of its evidence.

The evaluation demonstrates that the trust algebra provides
all five dimensions of provenance metadata---an 8-level lattice,
algebraic composition, an append-only audit log, copilot
channel ceilings, and the no-silent-promotion invariant---that
no existing proof assistant or LLM system tracks.
In particular, the integration of AI-generated
evidence at a controlled, non-promotable trust level
addresses a pressing need as LLMs become standard
components of the verification tool-chain.

Future work will explore three directions.
First, we plan to develop a \emph{trust type system}
that statically rejects programs containing silent
promotions, complementing the current dynamic enforcement.
Second, we will investigate \emph{quantitative trust}
by replacing the discrete lattice with a continuous
semiring, enabling probabilistic reasoning about
evidence quality.
Third, we will extend the evaluation to industrial
codebases with millions of lines of code, testing
the scalability of the audit log and the ceiling
enforcement mechanism.


%% ═══════════════════════════════════════════════════════════════════════
%% Bibliography
%% ═══════════════════════════════════════════════════════════════════════

\begin{thebibliography}{99}

\bibitem{martinlof84}
P.~Martin-L\"of.
\newblock \emph{Intuitionistic Type Theory}.
\newblock Bibliopolis, Naples, 1984.

\bibitem{coquand88}
T.~Coquand and C.~Huet.
\newblock The {Calculus of Constructions}.
\newblock \emph{Information and Computation}, 76(2--3):95--120, 1988.

\bibitem{gentzen35}
G.~Gentzen.
\newblock Untersuchungen \"uber das logische {S}chlie{\ss}en.
\newblock \emph{Mathematische Zeitschrift}, 39:176--210, 405--431, 1935.

\bibitem{lean4}
L.~de~Moura and S.~Ullrich.
\newblock The {Lean~4} theorem prover and programming language.
\newblock In \emph{CADE-28}, LNCS 12699, pp.~625--635. Springer, 2021.

\bibitem{coq}
The Coq Development Team.
\newblock \emph{The {Coq} Reference Manual}, version 8.18, 2023.

\bibitem{fstar}
N.~Swamy, C.~Hri\c{t}cu, C.~Keller, A.~Rastogi, et~al.
\newblock Dependent types and multi-monadic effects in {F*}.
\newblock In \emph{POPL 2016}, pp.~256--270. ACM, 2016.

\bibitem{dafny}
K.~R.~M. Leino.
\newblock Dafny: An automatic program verifier for functional correctness.
\newblock In \emph{LPAR-16}, LNCS 6355, pp.~348--370. Springer, 2010.

\bibitem{hoare69}
C.~A.~R. Hoare.
\newblock An axiomatic basis for computer programming.
\newblock \emph{Communications of the ACM}, 12(10):576--580, 1969.

\bibitem{reynolds02}
J.~C. Reynolds.
\newblock Separation logic: A logic for shared mutable data structures.
\newblock In \emph{LICS 2002}, pp.~55--74. IEEE, 2002.

\bibitem{iris}
R.~Jung, R.~Krebbers, J.-H.~Jourdan, A.~Bizjak, L.~Birkedal, and D.~Dreyer.
\newblock Iris from the ground up: A modular foundation for higher-order
  concurrent separation logic.
\newblock \emph{Journal of Functional Programming}, 28:e20, 2018.

\bibitem{isabelle}
T.~Nipkow, L.~C. Paulson, and M.~Wenzel.
\newblock \emph{Isabelle/HOL --- A Proof Assistant for Higher-Order Logic}.
\newblock LNCS 2283. Springer, 2002.

\bibitem{liquidhaskell}
N.~Vazou, E.~L. Seidel, R.~Jhala, D.~Vytiniotis, and S.~Peyton~Jones.
\newblock Refinement types for {Haskell}.
\newblock In \emph{ICFP 2014}, pp.~269--282. ACM, 2014.

\bibitem{sel4}
G.~Klein, K.~Elphinstone, G.~Heiser, et~al.
\newblock {seL4}: Formal verification of an {OS} kernel.
\newblock In \emph{SOSP 2009}, pp.~207--220. ACM, 2009.

\bibitem{certikos}
R.~Gu, Z.~Shao, H.~Chen, X.~Wu, J.~Kim, V.~Sj\"oberg, and D.~Costanzo.
\newblock {CertiKOS}: An extensible architecture for building certified
  concurrent {OS} kernels.
\newblock In \emph{OSDI 2016}, pp.~653--669. USENIX, 2016.

\bibitem{buneman01}
P.~Buneman, S.~Khanna, and W.-C. Tan.
\newblock Why and where: A characterization of data provenance.
\newblock In \emph{ICDT 2001}, LNCS 1973, pp.~316--330. Springer, 2001.

\bibitem{green07}
T.~J. Green, G.~Karvounarakis, and V.~Tannen.
\newblock Provenance semirings.
\newblock In \emph{PODS 2007}, pp.~31--40. ACM, 2007.

\bibitem{halpern03}
J.~Y. Halpern.
\newblock \emph{Reasoning about Uncertainty}.
\newblock MIT Press, 2003.

\bibitem{shafer76}
G.~Shafer.
\newblock \emph{A Mathematical Theory of Evidence}.
\newblock Princeton University Press, 1976.

\bibitem{paper02}
H.~Young.
\newblock The judgment 8-tuple: An algebraic foundation for evidence-carrying
  proofs.
\newblock \emph{Judgment Geometry Papers}, Paper~02, 2025.

\bibitem{gpt4o}
OpenAI.
\newblock {GPT-4o} system card.
\newblock Technical report, OpenAI, 2024.

\end{thebibliography}

\end{document}
